\begin{table}[H]
\centering
\captionsetup{justification=centering}
\caption{Hiperparâmetros estimados para o total de desocupados - modelos univariado e multivariado}
\label{tab:hiperdesoc}
\scalebox{0.92}{
\renewcommand{\arraystretch}{0.85}
\begin{tabular}{llccccc}
\toprule
& & \multicolumn{5}{c}{\textbf{Modelo Univariado}} \\
\cmidrule(lr){3-7}
\textbf{Estrato geográfico} & \textbf{Processo} & \(\hat{\sigma}_L^2\) & \(\hat{\sigma}_R^2\) & \(\hat{\sigma}_S^2\) & \(\hat{\sigma}_I^2\) & \(\hat{\sigma}_{\tilde{e}}^2\) \\
\midrule
01 - Belo Horizonte & AR(1) & 0,0483 & 18,4914 & 0,0000 & 0,0000 & 0,9052 \\
02 - Entorno e Colar Metropolitano de BH & MA(1) & 0,0161 & 102,7494 & 0,0000 & 0,0000 & 0,8524 \\
03 - Sul de Minas & AR(1) & 82,4451 & 1,1848 & 0,0000 & 0,0000 & 0,8467 \\
04 - Triângulo Mineiro & AR(1) & 0,0246 & 8,3891 & 0,0000 & 0,0000 & 0,9517 \\
05 - Zona da Mata & AR(1) & 46,8067 & 0,5211 & 0,0000 & 0,0000 & 0,9046 \\
06 - Norte de Minas & AR(1) & 127,4207 & 1,4080 & 0,0000 & 0,0000 & 0,7976 \\
07 - Vale do Rio Doce & AR(1) & 112,0160 & 1,1011 & 0,0000 & 0,0000 & 0,6863 \\
08 - Central & AR(3) & 0,0306 & 21,8547 & 0,0000 & 0,0000 & 0,9062 \\
\midrule
& & \multicolumn{5}{c}{\textbf{Modelo Multivariado}} \\
\cmidrule(lr){3-7}
\textbf{Estrato geográfico} & \textbf{Processo} & \(\hat{\sigma}_L^2\) & \(\hat{\sigma}_R^2\) & \(\hat{\sigma}_S^2\) & \(\hat{\sigma}_I^2\) & \(\hat{\sigma}_{\tilde{e}}^2\) \\
\midrule
01 - Belo Horizonte & AR(1) & 0,0239 & 33,8915 & 0,0000 & 0,0000 & 0,9052 \\
02 - Entorno e Colar Metropolitano de BH & MA(1) & 0,0060 & 112,5743 & 0,0000 & 0,0000 & 0,8524 \\
03 - Sul de Minas & AR(1) & 0,0658 & 19,8207 & 0,0000 & 0,0000 & 0,8467 \\
04 - Triângulo Mineiro & AR(1) & 0,0383 & 13,5033 & 0,0000 & 0,0000 & 0,9517 \\
05 - Zona da Mata & AR(1) & 1,1664 & 10,7285 & 0,0000 & 0,0000 & 0,9046 \\
06 - Norte de Minas & AR(1) & 0,5315 & 27,4822 & 0,0000 & 0,0000 & 0,7976 \\
07 - Vale do Rio Doce & AR(1) & 0,7853 & 30,7076 & 0,0000 & 0,0000 & 0,6863 \\
08 - Central & AR(3) & 0,0474 & 15,3830 & 0,0000 & 0,0000 & 0,9062 \\
\bottomrule
\end{tabular}}
\fonte{Elaboração própria, com base nos dados da PNAD Contínua. Nota: o processo do erro amostral é identificado individualmente por estrato (Seção \ref{sec:metodologia}), e \(\hat{\sigma}_{\tilde{e}}^2\) é derivada dos coeficientes do processo, por isso coincidindo nos dois modelos. Sobre \(\sigma_I^2\), ver a ressalva de identificação no texto.}
\end{table}
